\scanpage{42}

\[
\boxed{\tau_h(Y)\sum(-1)^i\varphi_Y(E^{(i)})=f_*(\tau_h(X)\cdot\varphi_X(E))}\tag{1ter}
\]

Prenant les \uline{augmentations} des deux mbs, la formule se réduit essentiellement à R.R. \uline{usuel} (\(Y\) réduit à un point !) sur la fibre. En effet, \unsure{la question étant locale}, le mb de gauche est égal à
\[
\sum(-1)^i\operatorname{rg} E^{(i)}=\chi(Z;E|Z)
\]
(\(Z\) étant une fibre de \(X\to Y\)), le mb de droite est égal aussi à
\[
\unsure{K_p|\kappa_p}\bigl(i^*[(\tau_h(X)\cdot\varphi_X(E))^{(p)}]\bigr),
\]
où \(i:Z\mathbin{\text{---}}X\) est l'application identique, et où \(p\) est la dimension de la fibre, et \({}^{(p)}\) signifie qu'on prend la composante de degré \(p\). On a
\[
\begin{aligned}
i^*[\quad]&=i^*(\tau_h(X))\deleted{\unsure{\varphi}}\cdot i^*(\varphi_X(E))\\
&=\tau_h(C_Z(i^*(T_X)))\varphi_Z(i^*(E))\\
&=\tau_h(C_Z(T_Z))\varphi_Z(E|Z)=\tau_h(Z)\varphi_Z(E|Z).
\end{aligned}
\]
On obtient bien \inserted{\unsure{exactement}} la formule de Hirzebruch pour la fibre. --- Mais prenant des \deleted{\unsure{ci-dessus}}\inserted{composantes} de degré \(\geq1\) des deux mbs, on détermine les classes de Chern (mod torsion) de \(\sum(-1)^iE^{(i)}\) (donc, en certains cas) de \(E^{(0)}\), ce qui n'est pas contenu dans la formule de Hirzebruch proprement dite. On trouve \unsure{en passant} le résultat suivant (qui n'est pas non plus trivial): \uline{pour \(n\) grand}, \uline{la classe de Chern modifiée de} \(E(n)^0=f_*(E(n))\) \uline{dépend polynomialement}

\margin{\unsure{peut aussi se déduire dans le cas général}}
\margin{41}
